\documentclass[11pt]{article}
\usepackage[utf8]{inputenc}
\usepackage[margin=1in]{geometry}
\usepackage{amsmath, amssymb, amsthm}
\usepackage{booktabs}
\usepackage{hyperref}

% --- Custom Header Information ---
\title{
    \Large \textbf{School of Engineering and Applied Science} \\
    \vspace{0.5cm}
    \large \textbf{Lecture Scribe}: Lecture 22 (CSE 400) \\
    \vspace{0.2cm}
    \normalsize \textbf{Name:} Manya Chudasama\\
    \vspace{0.2cm}
    \normalsize \textbf{Enrollment No:} AU2440013
}

\date{} 
\begin{document}
\maketitle
\section*{1) List of Topics Covered}
\begin{enumerate}
    \item \textbf{Chebyshev’s Inequality}
    \begin{itemize}
        \item Intuition and Key Idea
        \item Proof using Markov’s Inequality
    \end{itemize}
    \item \textbf{Running Example: Coin Flipping}
    \begin{itemize}
        \item Setup and Binomial Distribution application
        \item Comparison with Markov’s Bound
    \end{itemize}
    \item \textbf{Real-Life Applications}
    \begin{itemize}
        \item Exam Scores Scenario
        \item Network Latency SLA
    \end{itemize}
\end{enumerate}

\section*{2) Explanation of Topics Covered}
\begin{enumerate}
    \item \textbf{Chebyshev’s Inequality} \\
    Chebyshev’s Inequality is a powerful tool in probability that uses both the mean ($\mu$) and variance ($Var(X)$) to control the deviation of a random variable from its mean. It helps bound the probability of ``tail regions''—values that are far from the average. Unlike Markov's inequality, which only requires the mean and applies only to non-negative variables, Chebyshev's bound often provides a tighter estimate because it incorporates the spread (variance) of the data.

    \item \textbf{Intuition} \\
    Geometrically, it limits how much area can exist in the tails of a distribution at a distance `$a$' from the mean. As the sample size or certain parameters increase, these bounds typically decrease, showing that values are more likely to stay near the mean. 
\end{enumerate}

\section*{3) List of Definitions and Theorems} 
\begin{enumerate}
    \item \textbf{Definition: Markov’s Inequality} \\
    For a non-negative random variable $X$ and any $a > 0$:
    \[ P(X \ge a) \le \frac{E[X]}{a} \]

    \item \textbf{Theorem Statement: Chebyshev’s Inequality} \\
    For a random variable $X$ with mean $\mu$ and variance $\sigma^2$, for any $a > 0$:
    \[ P(|X - \mu| \ge a) \le \frac{Var(X)}{a^2} \]

    \item \textbf{Detailed Step-by-Step Proof}
    \begin{enumerate}
        \item[1.] Consider the event $|X - \mu| \ge a$.
        \item[2.] Square both sides to relate it to variance: $(X - \mu)^2 \ge a^2$.
        \item[3.] Apply \textbf{Markov’s Inequality} to the non-negative random variable $(X - \mu)^2$:
            \[ P((X - \mu)^2 \ge a^2) \le \frac{E[(X - \mu)^2]}{a^2} \]
        \item[4.] Substitute the definition of variance $Var(X) = E[(X - \mu)^2]$:
            \[ P(|X - \mu| \ge a) \le \frac{Var(X)}{a^2} \]
    \end{enumerate}
\end{enumerate}

\section*{4) Important Examples}

\subsection*{Example 1: Exam Grades with Variance} 
\vspace{0.1 cm}
\textbf{Problem} : The average grade is 70\% and the standard deviation is 5\%. The ``A'' cutoff is 90\%. Apply Chebyshev’s inequality to obtain a tighter upper bound on the fraction of ``A'' students.

\vspace{0.2cm}
\noindent \textbf{Solution:} 
\begin{itemize}
    \item \textbf{Step 1:} Identify mean and variance. \\ $\mu = 70$, $\sigma = 5 \Rightarrow \sigma^2 = 25$.
    \item \textbf{Step 2:} Define the event. \\ An ``A'' student has $X \ge 90$. Note that $X \ge 90 \Rightarrow |X - 70| \ge 20$. We use the two-sided Chebyshev bound to cover this one-sided event.
    \item \textbf{Step 3:} State Chebyshev with $a = 20$. \\ $P(|X - \mu| \ge a) \le \frac{\sigma^2}{a^2}$.
    \item \textbf{Step 4:} Substitute values. \\ $P(X \ge 90) \le P(|X - 70| \ge 20) \le \frac{25}{20^2} = \frac{25}{400}$.
    \item \textbf{Result:} $P(X \ge 90) \le 0.0625$ (or $\frac{1}{16}$). This is much tighter than the bound provided by Markov’s inequality (0.778).
\end{itemize}

\vspace{0.4cm}

\subsection*{Example 2: Network Latency SLA} 

\textbf{Problem :} A network route has average latency $E[X] = 50\text{ ms}$ and variance $Var(X) = 16\text{ ms}^2$. Bound the probability that a packet's latency falls \textbf{outside} the range $[30\text{ ms}, 70\text{ ms}]$.

\vspace{0.2cm}
\noindent \textbf{Solution:} 
\begin{itemize}
    \item \textbf{Step 1:} Identify given parameters. \\ $\mu = 50$, $\sigma^2 = 16$.
    \item \textbf{Step 2:} Rewrite failure condition as absolute deviation. \\ $X \notin [30, 70] \iff |X - 50| \ge 20$.
    \item \textbf{Step 3:} Apply Chebyshev with $a = 20$. \\ $P(|X - 50| \ge 20) \le \frac{16}{20^2} = \frac{16}{400}$.
    \item \textbf{Result:} $P(X \text{ is outside range}) \le 0.04$.
\end{itemize}

\section*{5) List of Important Formulas}
\begin{table}[h!]
\centering
\renewcommand{\arraystretch}{1.5}
\setlength{\tabcolsep}{12pt}
\begin{tabular}{ll}
\toprule
\textbf{Formula / Concept} & \textbf{Result} \\
\midrule
\textbf{Markov’s Inequality} & $P(X \ge a) \le \frac{E[X]}{a}$ \\
\textbf{Chebyshev’s Inequality} & $P(|X - \mu| \ge a) \le \frac{Var(X)}{a^2}$ \\
\textbf{Binomial Mean (Fair Coin)} & $\mu = \frac{n}{2}$ \\
\textbf{Binomial Variance (Fair Coin)} & $Var(X) = \frac{n}{4}$ \\
\bottomrule
\end{tabular}
\end{table}

\end{document}
